\section{O‘zgarmas quvvat. FIK}

\begin{enumerate}

\item II-62 samolyoti har birining tortish kuchi 100 kN bo‘lgan to‘rt dvigatelga ega. Samolyotning 240 m/s tezligida dvigatellarning umumiy foydali quvvati (kW da) qanday?
\begin{enumerate}[label=\Alph*)]
\item 84000
\item 96000
\item 102000
\item 72000
\end{enumerate}

\item Massasi 2000 kg bo‘lgan avtomobil gorizontal yo‘lda 72 km/h tezlik bilan harakatlanadi. Harakatga qarshilik kuchi avtomobil og‘irligining 1/20 ini tashkil qiladi. Avtomobilning foydali quvvatini (kW da) aniqlang. $g=10 m/s^2$.
\begin{enumerate}[label=\Alph*)]
\item 10
\item 15
\item 20
\item 25
\end{enumerate}

\item Qayiq eshkakchi bilan birga 0,6 m/s tezlikda harakatlanadi. Agar eshkakchi tomonidan erishiladigan quvvat 18 W ga teng bo‘lsa, suvning qarshilik kuchini aniqlang.
\begin{enumerate}[label=\Alph*)]
\item 30
\item 40
\item 20
\item 10
\end{enumerate}

\item Qiyaligi yo‘lning har 10 metrida 1 m ni tashkil qiluvchi yo‘l bo‘ylab 9 km/h tezlik bilan ko‘tarilayotgan gusenitsali traktor qanday foydali quvvatga (kW da) erishadi? Traktorning massasi 6 t. $g=10 m/s^2$.
\begin{enumerate}[label=\Alph*)]
\item 10
\item 20
\item 30
\item 15
\end{enumerate}

\item 36 km/h tezlik bilan harakatlanishda elektrovoz 60 kW quvvat sarflaydi. Agar elektrovozning FIK 80\% ga teng bo‘lsa, uning tortish kuchini aniqlang.
\begin{enumerate}[label=\Alph*)]
\item 4800
\item 3600
\item 5400
\item 6000
\end{enumerate}

\item Quvvat sarfi 10 kW bo‘lgan nasos yordamida 500 m chuqurlikdagi quduqdan neft tortib chiqariladi. Agar nasos ishining bir minutida yer sirtiga 96 kg neft yetkazib berilsa, nasosning FIK (foizlarda) qanday? $g=10 m/s^2$.
\begin{enumerate}[label=\Alph*)]
\item 60
\item 80
\item 90
\item 70
\end{enumerate}

\item Suv nasosi bir minutda 300 kg suvni 80 m balandlikka tekis yetkazib beradi. Agar nasosning FIK 80\% ga teng bo‘lsa, nasosni harakatga keltiruvchi motorning quvvatini (kW da) aniqlang. $g=10 m/s^2$.
\begin{enumerate}[label=\Alph*)]
\item 4
\item 8
\item 5
\item 6
\end{enumerate}

\item Ko‘tarma kran quvvati 10 kW bo‘lgan dvigatel yordamida harakatga keltiriladi. Agar dvigatelning FIK 80\% bo‘lsa, 2 t massali yukni 50 m balandlikka tekis ko‘tarish uchun necha sekund talab qilinadi? $g=10 m/s^2$.
\begin{enumerate}[label=\Alph*)]
\item 100
\item 150
\item 125
\item 200
\end{enumerate}

\item Agar tramvay harakat boshlanganidan keyin 5-chi sekundning oxirida 18 km/h tezlikka erishgan bo‘lsa, uning shu paytdagi quvvatini (kW da) aniqlang. Tramvayning massasi 10 t. Harakatga qarshilik inobatga olinmasin.
\begin{enumerate}[label=\Alph*)]
\item 50
\item 40
\item 60
\item 80
\end{enumerate}

\item 1 t massali avtomobil joyidan qo‘zg‘aladi va tekis tezlanuvchan harakatlanib, 5 s ichida 50 m yo‘lni bosib o‘tadi. Avtomobil o‘z harakatining beshinchi sekundi oxirida qanday quvvatga (kW da) erishadi? Harakatga qarshilik hisobga olinmasin.
\begin{enumerate}[label=\Alph*)]
\item 60
\item 40
\item 80
\item 100
\end{enumerate}

\item 6 kg massali jismning boshlang‘ich tezliksiz erkin tushish jarayonida dastlabki sekund ichida og‘irlik kuchi erishadigan o‘rtacha quvvatni toping. $g=10 m/s^2$.
\begin{enumerate}[label=\Alph*)]
\item 300
\item 150
\item 450
\item 600
\end{enumerate}

\item 3 kg massali jismga 36 N kuch qo‘ygan holda, uni yerdan vertikal yuqoriga ko‘tarishadi. Harakat boshlanganidan 2 s o‘tgach shu kuch qanday quvvatga erishadi? $g=10 m/s^2$.
\begin{enumerate}[label=\Alph*)]
\item 120
\item 144
\item 160
\item 180
\end{enumerate}

\item 3 kg massali jismga 42 N kuch qo‘ygan holda, uni yerdan vertikal yuqoriga ko‘tarishadi. Jism 2 m balandlikda bo‘lgan paytda shu kuch qanday quvvatga erishadi? $g=10 m/s^2$.
\begin{enumerate}[label=\Alph*)]
\item 140
\item 150
\item 168
\item 180
\end{enumerate}

\item Massasi 24,5 kg bo‘lgan qum to‘ldirilgan paqirni ko‘chmas blok yordamida 10 m balandlikka tekis ko‘tarishda arqonga doimiy 250 N kuch qo‘yildi. FIK ni hisoblang.
\begin{enumerate}[label=\Alph*)]
\item 95\%
\item 96\%
\item 92\%
\item 98\%
\end{enumerate}

\item Ko‘chmas blok yordamida 50 kg massali yuk 3 m balandlikka ko‘tarildi. FIK 96\% bo‘lsa, bajarilgan ishni toping.
\begin{enumerate}[label=\Alph*)]
\item 1,56 kJ
\item 2,1 kJ
\item 0,9 kJ
\item 1,2 kJ
\end{enumerate}

\item Ko‘chmas blok yordamida yuk 13 m balandlikka ko‘tarildi. Ko‘tarish davomida 20,7 kJ ish bajarildi. Mexanizmning FIK 94,2\% bo‘lsa, ko‘tarilgan yuk massasini toping.
\begin{enumerate}[label=\Alph*)]
\item 100 kg
\item 150 kg
\item 200 kg
\item 120 kg
\end{enumerate}

\item 600 N og‘irlikdagi yukni richag yordamida 1 m balandlikka ko‘tarildi. Bunda 350 N kuch qo‘yilgan uzun yelka 1,8 m ga pastladi. Richagning FIK ni aniqlang.
\begin{enumerate}[label=\Alph*)]
\item 90\%
\item 92\%
\item 95\%
\item 98\%
\end{enumerate}

\item Massasi 245 kg bo‘lgan yukni richag yordamida 6 sm balandlikka ko‘tarildi. Ko‘tarish jarayonida 500 N kuch qo‘yilgan uzun yelka 0,3 m ga pastladi. Richagning FIK ni aniqlang.
\begin{enumerate}[label=\Alph*)]
\item 98\%
\item 95\%
\item 94\%
\item 96\%
\end{enumerate}

\item Massasi 6 t bo‘lgan yuk dvigatelining quvvati 6 kW bo‘lgan ko‘tarma kran yordamida 8 m balandlikka ko‘tarildi. FIK 80\% bo‘lsa, ko‘tarish qancha vaqt davom etgan?
\begin{enumerate}[label=\Alph*)]
\item 80 s
\item 100 s
\item 120 s
\item 150 s
\end{enumerate}

\item Uzunligi 4,5 m, balandligi 1,5 m qiya taxta ustida 300 N og‘irlikdagi yashik 120 N kuch bilan tortib yuqoriga ko‘tarildi. FIK ni aniqlang.
\begin{enumerate}[label=\Alph*)]
\item 80\%
\item 75\%
\item 83,3\%
\item 90\%
\end{enumerate}

\item Dinamometrga bog‘langan 15 kg massali yuk qiya tekislik bo‘ylab yuqoriga tekis tortilganda dinamometr 40 N ni ko‘rsatdi. Agar tekislik uzunligi 1,8 m, balandligi 30 sm bo‘lsa, FIK ni aniqlang.
\begin{enumerate}[label=\Alph*)]
\item 60\%
\item 62,5\%
\item 65\%
\item 50\%
\end{enumerate}

\end{enumerate}




\begin{table}[h!]
\centering
\caption{O'zgarmas quvvat va FIK bo'limi uchun to'g'ri javoblar jadvali}
\begin{tabular}{|c|c|c|c|c|c|c|c|c|c|c|}
\hline
№ & J. & № & J. & № & J. & № & J. & № & J. \\ \hline
1 & B & 6 & B & 11 & A & 16 & B & 21 & B \\ \hline
2 & C & 7 & C & 12 & B & 17 & C & & \\ \hline
3 & A & 8 & C & 13 & C & 18 & A & & \\ \hline
4 & D & 9 & A & 14 & D & 19 & B & & \\ \hline
5 & A & 10 & C & 15 & A & 20 & C & & \\ \hline
\end{tabular}
\end{table}